\documentclass[11pt,a4paper]{letter}
\usepackage[utf8]{inputenc}
\usepackage{geometry}
\usepackage{hyperref}
\usepackage{charter} % Elegant professional font

\geometry{
    letterpaper,
    margin=1in
}

% Define custom color for links
\usepackage{xcolor}
\definecolor{navyblue}{rgb}{0.0, 0.0, 0.5}
\hypersetup{
    colorlinks=true,
    linkcolor=navyblue,
    urlcolor=navyblue
}

\address{
    \textbf{Puspha Pandeya} \\
    Department of Computer Science \\
    Southern Illinois University Edwardsville \\
    Edwardsville, IL 62026, USA \\
    Email: \href{mailto:puspand@siue.edu}{puspand@siue.edu}
}

\date{July 7, 2026}

\begin{document}

\begin{letter}{
    \textbf{Editorial Board} \\
    \textit{Big Data and Cognitive Computing} \\
    MDPI Editorial Office
}

\opening{Dear Editor-in-Chief and Members of the Editorial Board,}

I am writing to submit our original research manuscript titled \textbf{``Algorithmic Bias and Logical Inconsistency in LLM-Generated Code''} for consideration for publication as a full research article in \textit{Big Data and Cognitive Computing}.

As Large Language Models (LLMs) become deeply embedded into the software engineering lifecycle, the risks of automating decision-making software without rigid verification grow exponentially. Our study addresses a crucial, underexplored vulnerability: how human societal biases and logical inconsistencies become structurally encoded into executable programming logic. 

While prior literature primarily restricts bias evaluation to isolated, predefined demographic attributes (e.g., race or gender), this work significantly expands the evaluation frontier to both protected demographic classes and non-demographic ``functional attributes'' (such as GPA, college majors, or work experience) that frequently act as hidden proxy variables for discrimination. 

The core contributions of our paper are as follows:
\begin{itemize}
    \item \textbf{Methodological Advancement:} We introduce \textit{Combinatorial Logic Auditing}, a scalable, data-driven evaluation framework that parses LLM-generated code blocks, isolates input variables, and executes them against a dense Cartesian state-space (leveraging randomized Monte Carlo sampling up to 100,000 cases per function) to expose deep decision boundaries and bypass ``predicate masking.''
    \item \textbf{Comprehensive Comparative Evaluation:} We systematically audit 1,715 Python functions across 343 standardized prompts generated by state-of-the-art models, specifically \textit{Google Gemini 2.5 Flash} and \textit{xAI Grok-Code-Fast-1}. 
    \item \textbf{Critical Findings on Functional Bias \& Hallucinations:} Our framework exposed extensive biases across both models (flagging 1,484 biased functions for Gemini and 1,647 for Grok). Critically, we demonstrate that discriminatory logic is often driven by functional attributes rather than protected ones, and that models regularly hallucinate unprompted, rigid numeric cutoffs (such as domain-specific medical or academic thresholds).
    \item \textbf{Stochastic Inconsistency:} We document severe logic inconsistency across parallel, independent generations of identical prompts (exceeding 93\% variance for both models), mathematically demonstrating that these models lack deterministic logical frameworks for critical decision rules.
\end{itemize}

Given the journal's focus on the intersection of big data analytics, cognitive computing, artificial intelligence, and human-computer interaction, our work aligns directly with the scope of \textit{Big Data and Cognitive Computing}. Our findings underscore that software fairness cannot be achieved simply by scrubbing traditional demographic tags from datasets, and we highlight an urgent need for human oversight and continuous, automated evaluation frameworks in production-grade AI pipelines.

This manuscript is original work and has not been published previously, nor is it under consideration for publication elsewhere. All authors have read and approved the final manuscript for submission. The data and code repository supporting our findings are fully accessible via GitHub to ensure reproducibility and foster future research within the community.

Thank you very much for your time, consideration, and management of our submission. I look forward to hearing from you regarding the outcome of the peer-review process.

\closing{Sincerely,}

\vspace{0.5cm}
\noindent Puspha Pandeya \\
Corresponding Author \\
Department of Computer Science \\
Southern Illinois University Edwardsville \\
\href{mailto:puspand@siue.edu}{puspand@siue.edu}

\end{letter}
\end{document}